\section{Drainage Design and Water-Table Management}

Drainage design and active water-table management are the primary levers
through which practitioners set both the volume of subsurface outflow and the
loads it carries. Because solute export in drained cropland is governed first
by hydrology, the geometric choices of spacing and depth and the operational
choice of when to hold or release water translate directly into edge-of-field
flow and nutrient loss. This section synthesises what the evidence has
established about those controls: how spacing and depth trade flow against
loss, how controlled drainage and drainage water management reduce loads in
proportion to the flow they retain, how capture-and-reuse systems recover both
water and quality benefits, and how enhanced configurations and controlled
drainage perform under projected climates.

\subsection{Drain Spacing and Depth as Controls on Flow and Loss}

Drain spacing and depth are well-established design controls on both drainage
volume and nitrate loss. More closely spaced drains produce greater drainage
and greater nitrate-N loss than wider spacings, whereas nitrate
\emph{concentrations} generally do not vary consistently with spacing---load
differences are driven by flow rather than by concentration
\parencite{Kladivko2021,Kladivko2004}. Calibrated models consistently show that
shallower and wider-spaced drains reduce tile flow and nitrate loss at modest
yield cost: reducing drain depth from $1.5$ to $0.9$~m lowered simulated
nitrate loss from $34.0$ to $29.4$~kg~ha$^{-1}$, and widening spacing from $27$
to $60$~m lowered it from $33.8$ to $28.4$~kg~ha$^{-1}$
\parencite{Moriasi2013,Nangia2010c,Davis2000,Hwang2024}. The design tension is
that reducing subsurface loss shifts water to surface runoff: for phosphorus,
increasing spacing reduced subsurface volumes and tile TP loss (more strongly
on finer soils) while increasing surface runoff and its associated TP load
\parencite{Morrison2013}.

These relationships motivate an optimum-intensity principle: use the minimum
subsurface intensity required to relieve excess-water stress, fit the system to
the soil, and apply water-table control where maximum intensity is unnecessary
\parencite{Breve1998}. Deeper drains often raise profit but substantially
increase nitrate losses, and simulation showed that widening spacing to $100$~m
cut total nitrate-N loss by $20\%$ ($37$ to $29$~kg~ha$^{-1}$) at small profit
cost \parencite{ElSadek2002b,Breve1998}; deep drains are frequently unnecessary
for salinity control \parencite{Ritzema2011}. Shallow drainage reproduces much
of the flow- and load-reduction benefit, reducing Iowa subsurface volume by
$46$--$58\%$ and nitrate by $29$--$49\%$ generally without reducing yield
\parencite{Helmers2012,Schott2017}. A multi-state analysis found shallow drains
gave quicker water-table drawdown, reduced yield variability, and $3.4$~cm less
total flow than deep drains, though with more surface runoff and a lower
benefit--cost ratio \parencite{Ghane2021b}, and field comparisons found shallow
systems had the lowest nitrate concentration and loading among configurations,
sometimes at a yield penalty whereas controlled drainage trended toward higher
yields \parencite{Smith2019,Smith2011b}.

\subsection{Controlled Drainage and Drainage Water Management}

Controlled drainage (CD)---raising an outlet control structure to hold the
water table between drainage events---is among the best-established
conservation-drainage practices. A synthesis found CD reduced drainage outflow
by $17\%$ to more than $80\%$ depending on soil, drainage intensity, control
strategy, and location, with nitrogen-loss reduction often about the same
percentage as the volume reduction \parencite{Skaggs2010}. A global
meta-analysis of 84 publications and 1788 paired site-years confirmed that CD
and controlled drainage with subirrigation (CDSI) reduce subsurface discharge
and nutrient losses while modestly increasing yield; CDSI reduced nitrate and
TN more, CD reduced TP more, and CDSI (but not CD) significantly increased
surface runoff and surface nitrate loss \parencite{Wu2025}. Simulations across
48 Midwest locations found CD reduced drainage and N losses everywhere, with
efficacy strongly climate-dependent \parencite{Youssef2018}.

Field studies quantify these effects concretely. Quebec CDSI cut drainage
volume $24\%$, nitrate concentration $25\%$, and nitrate loss $43\%$
\parencite{Ng2000}; a long-term Canadian study found CD reduced total nitrate
loss by $\sim58\%$ and increased yield $3\%$ without elevated greenhouse-gas
emissions \parencite{Madramootoo2024}; and nine years of North Carolina CD
reduced outflow $33\%$ and nitrate export $30\%$ while raising yields
$10$--$11\%$ \parencite{Poole2018}. Growing-season CD reduced water flux
$\sim60\%$, nitrate $51\%$, ammonium $58\%$, DRP and TP $66\%$, and
\emph{E.\ coli} $76\%$, with watershed-scale deployment cutting stream fluxes
of water, N, and P \parencite{Sunohara2016,Sunohara2015}. In Illinois, managed
systems always exported less nitrate than free drainage, with flow reductions
of $35$--$96\%$ and load reductions of $37$--$79\%$ (averaging $61\%$)
\parencite{Cooke2012}.

The mechanism is consistently one of flow--load coupling: syntheses report
drainage water management reducing tile flow by $\sim46\%$ and nitrate loads by
$\sim45$--$48\%$ \emph{mainly by reducing discharge rather than by enhancing
denitrification}, which plays only a secondary role
\parencite{Ross2016,Ward2018}. Two trade-offs qualify the benefit. First,
because reduction operates through volume, drainage management often fails to
lower nitrate \emph{concentrations} even where it lowers loads
\parencite{Grenon2023,Wesström2014}. Second, CD can increase surface DRP and TP
that partly offset subsurface benefits \parencite{King2022,Wu2025}. Timing is
decisive: DRAINMOD analysis identified early March as the most effective
control window (reductions up to $100\%$) whereas April starts resembled free
drainage \parencite{Kesicka2023c}. Greenhouse-gas responses are mixed but
generally not adverse \parencite{Madramootoo2024,Smith2019}.

\subsection{Drainage Water Recycling and Subirrigation}

Drainage water recycling (DWR)---capturing and storing drainage for
supplemental irrigation---benefits both production and downstream water
quality. Yield gains reach $61\%$ for corn and $31\%$ for soybean, with
Canadian N-loss reductions of $40$--$70\%$ and P-loss reductions of
$12$--$36\%$, and modelled downstream N-delivery reductions of $24$--$37\%$
\parencite{Hay2021}. A North Carolina reservoir reduced TN load $47\%$, TP
$30\%$, and sediment $87\%$, with hydraulic retention time the key control,
although the reservoir was a net source of particulate P
\parencite{Moursi2023,Moursi2022}. Resilience is a defining advantage:
multi-location synthesis of 53 site-years showed DWR is more resilient than
free drainage, protecting yield during both drought and excess precipitation
and reducing yield variability by $22\%$ \parencite{Willison2021}. Free
drainage and drainage water recycling raised corn yield $20\%$ and $26\%$
respectively over a nondrained control (recycling gains reaching $48\%$ in
individual years) by avoiding drought losses \parencite{Kaur2021}, and North
Carolina recycling
raised corn and soybean by $8\%$ and $17\%$ with subirrigation raising the
water table $15$~cm \parencite{Moursi2024}. Adoption economics indicate crop
insurance is not a major barrier \parencite{Sellars2022}, and in arid systems
integrated on-farm drainage management reused all but $0.7\%$ of applied water
\parencite{Ayars2014,Kosari2024}.

Subirrigation---delivering water upward through the drains---similarly lowers
losses and stabilises supply. It reduced soil-profile nitrate by up to $50\%$
and drainage nitrate concentrations by $55$--$73\%$, with water-quality
benefits accruing even when outflow occasionally exceeded free drainage,
implicating enhanced denitrification \parencite{Elmi2004,Elmi2002,Noory2009}.
CDSI raised groundwater levels, supplied $170$--$410$~mm of infiltration, and
increased Dutch yields by up to $64\%$
\parencite{deWit2022,deWit2024,FAUSEY1995,Almen2021,Madramootoo2007}. Its
performance, however, is strongly site-dependent, governed by lateral hydraulic
conductivity: across three Ohio sites subirrigation improved relative yield to
$87$--$93\%$ and $80$--$91\%$ at two sites but had almost no effect at Defiance
owing to very low conductivity \parencite{Gunn2018,Gunn2019}. Where hydraulics
are favourable, CD and subirrigation kept soil water optimal in both wet and
dry years \parencite{Niaghi2019}, and in claypan soybean systems subirrigation
improved yield, reduced interannual variability, and allowed earlier planting
\parencite{Nelson2011}; combining drainage water management with foliar
pyraclostrobin produced synergistic soybean gains up to $1070$~kg~ha$^{-1}$
\parencite{Nelson2011b}.

\subsection{Enhanced Configurations and Persistence under Climate Change}

Beyond conventional single-level layouts, engineered configurations extend the
same design principles. Multi-level (bilevel) systems cut the time the water
table exceeded $0.5$, $0.75$, and $1.0$~m by $84\%$, $67\%$, and $41\%$ while
producing much lower shallow-drain salinity \parencite{Hornbuckle2007};
analytical theory shows that adding shallow drains between deep drains raises
the drainage rate and allows wider deep-drain spacing
\parencite{Hornbuckle2012}, with managed deep drains occupying an intermediate
position \parencite{Ayars2006b,Christen2001b}. Engineered enhancements alter
flow further: a permeable filter increased discharge (about $1.9$ times
conventional) and accelerated drawdown while reducing runoff
\parencite{Tao2016,Tao2017}; combined open-ditch-plus-subsurface designs
improved waterlogging removal \parencite{Ren2021,Tao2025}; siphon-type drainage
raised flow rates $27$--$90\%$ \parencite{Ren2023}; artificial macropore tubes
generated flow where none occurred and desalinated the profile
\parencite{Chen2021}; and supplementary drainage of a clay field raised mean
annual discharge $2.6$-fold and cut topsoil runoff by $75\%$
\parencite{Häggblom2019}.

A robust and management-relevant finding is that the benefits of controlled
drainage persist, and in some seasons strengthen, under projected climates. In
Ohio, conventional discharge was projected to decline while CD annual
reductions remained near $4$--$6\%$ and summer growing-season reductions
increased, indicating greater warm-season value \parencite{Pease2017}. Michigan
simulations projected free-drainage discharge rising $20\%$ by mid-century while
CD continued to reduce discharge by $63$--$72\%$, with future and historical
performance statistically indistinguishable \parencite{BinShokrana2023}.
Finnish scenarios projected CD effects remaining similar to the historical
interval \parencite{Salla2022}, and Canadian projections of warmer, wetter
conditions found CD reducing subsurface drainage more strongly than free
drainage, though rising temperature drove dry-stress yield reductions that
motivate the capture, storage, and reuse of drainage water
\parencite{Ndulue2024,Gunn2018}. Chinese and Polish studies likewise found that
appropriate outlet-blocking strategies could reduce future runoff and maintain
groundwater at present levels \parencite{Awad2021,Sojka2020}. Together these
results indicate that the flow--load coupling underpinning water-table
management is durable: because controlled systems act on discharge, their
proportional benefits carry forward into a wetter, warmer future.
